% paper/sections/appendix_numerics_solver_s5.tex
%
% 付録 D.1.3：欠陥補正法の収束理論
%   — 反復行列 M = I - ω L_L^{-1} L_H のスペクトル解析
%   — 変密度場における固有値分布
%   — 緩和係数 ω の理論的導出
%   — §8.3 (sec:defect_correction_main) の理論的裏付け

\subsection{収束理論：反復行列のスペクトル解析}
\label{app:dc_convergence_theory}

本付録は §\ref{sec:defect_correction_main} で定義した欠陥補正法（式~\eqref{eq:dc_three_step}）の
収束性を反復行列のスペクトル解析に基づいて理論的に保証する．

\subsubsection*{反復行列のスペクトル半径}

欠陥補正法の更新式（式~\eqref{eq:dc_three_step}）を誤差 $e^{(k)} \equiv p^{(k)} - p^*$
（$p^* = L_H^{-1} b$ は真の解）について整理すると：
\begin{equation}
  e^{(k+1)} = \underbrace{\bigl(I - \omega\,L_L^{-1}\,L_H\bigr)}_{\displaystyle \equiv M}\,e^{(k)}
  \label{eq:dc_iteration_matrix}
\end{equation}
反復の収束条件はスペクトル半径 $\rho(M) < 1$ であり，
$k$ 回反復後の誤差は $\|e^{(k)}\| \leq \rho(M)^k\,\|e^{(0)}\|$ で上界される．

\subsubsection*{\texorpdfstring{$M$ の固有値構造}{M の固有値構造}}

$L_L$ と $L_H$ が同一の固有ベクトル系 $\{v_j\}$ を共有すると仮定する
（等方一様格子で近似的に成立）．
$L_L\,v_j = \mu_j\,v_j$，$L_H\,v_j = \lambda_j\,v_j$ とすると
$M$ の固有値は：
\begin{equation}
  \sigma_j = 1 - \omega\,\frac{\lambda_j}{\mu_j}
  \label{eq:dc_eigenvalue}
\end{equation}
ここで $\lambda_j / \mu_j$ は CCD 演算子と FD 演算子の固有値比であり，
低周波モード（$\xi \to 0$）では $\lambda_j / \mu_j \to 1$（両者とも $\Ord{\xi^2}$ で一致），
高周波モード（$\xi \to \pi$）では $\lambda_j / \mu_j \neq 1$ となる（高次精度項の差異）．

\subsubsection*{固有値比の解析（完全連立 Eq-I+II 系）}

CCD は Eq-I（$p'$ 結合）と Eq-II（$p''$ 結合）を連立して解く．
$L_H$ の有効修正波数は両方程式を同時に解いた数値的固有値として定まり，
Eq-II 単体の近似式（$\Lambda_H \propto 24\sin^2(\xi/2)$）とは異なる．

滑らかモード（$\xi \to 0$）では CCD も FD もともに2階微分の厳密値に漸近するため
$\lambda_j/\mu_j \to 1$ となる．Nyquist モード（$\xi = \pi$）では完全連立解析
（Eq-I + Eq-II の行列固有値問題）により $\lambda_j/\mu_j = 2.4$ が得られる
（均一密度1次元，$N=32$ での数値検証と一致）．

\begin{center}
\PaperTableSetup
\begin{tabular}{ccccc}
\toprule
$kh/\pi$ & $\lambda_H h^2$ & $\lambda_{FD} h^2$ & 比 $r(k)$ & $\omega_{\max}$ \\
\midrule
$0\phantom{.00}$ (smooth) & $\approx -k^2 h^2$ & $\approx -k^2 h^2$ & $1.000$ & $2.000$ \\
$0.50$ & $-4.957$ & $-4.000$ & $1.239$ & $1.614$ \\
$0.75$ & $-11.660$ & $-6.828$ & $1.707$ & $1.171$ \\
$0.90$ & $-17.269$ & $-7.804$ & $2.213$ & $0.904$ \\
$1.00$ (Nyquist) & $-19.200$ & $-8.000$ & $2.400$ & $0.833$ \\
\bottomrule
\end{tabular}
\end{center}

固有値比は $\xi$ とともに\textbf{単調増加}し，
$[\alpha_{\min}, \alpha_{\max}] = [1.0, 2.4]$ の範囲をとる
（Nyquist が最大，滑らかモードが最小）．
\label{eq:dc_eigenvalue_ratio}

\subsubsection*{\texorpdfstring{最適緩和係数 $\omega^*$ の導出}{最適緩和係数 ω* の導出}}

$|\sigma_j| = |1 - \omega\,\lambda_j/\mu_j| < 1$ が全モードで成立するために，
$\omega < 2 / \alpha_{\max}$ が必要十分条件である：
\begin{equation}
  \omega < \frac{2}{\alpha_{\max}} = \frac{2}{2.4} \approx 0.833
  \label{eq:dc_omega_max}
\end{equation}
$\rho(M)$ を最小化する最適 $\omega^*$ は Chebyshev 中点公式より：
\begin{equation}
  \omega^* = \frac{2}{\alpha_{\min} + \alpha_{\max}}
  = \frac{2}{1.0 + 2.4}
  = \frac{2}{3.4}
  \approx 0.588
  \label{eq:dc_omega_opt}
\end{equation}
このとき
\begin{equation}
  \rho(M^*) = \frac{\alpha_{\max} - \alpha_{\min}}{\alpha_{\max} + \alpha_{\min}}
  = \frac{2.4 - 1.0}{2.4 + 1.0}
  = \frac{1.4}{3.4}
  \approx 0.41
  \label{eq:dc_rho_opt}
\end{equation}

実用的な $\omega$ 選択と各モードの減衰率を以下に示す：

\begin{center}
\PaperTableSetup
\begin{tabular}{cccc}
\toprule
$\omega$ & Nyquist $|\mu|$ & 滑らか $|\mu|$ & 評価 \\
\midrule
$1.0$（緩和なし） & $1.40$ → 発散 & $0.00$ & 不可 \\
$0.833$（$\omega_{\max}$） & $1.00$ → 停滞 & $0.17$ & 条件付 \\
$0.70$ & $0.28$ & $0.30$ & 合格（速い） \\
$\mathbf{0.50}$ & $\mathbf{0.20}$ & $\mathbf{0.50}$ & 合格（バランス型） \\
$0.30$ & $0.28$ & $0.70$ & 合格（安全側） \\
\bottomrule
\end{tabular}
\end{center}

exp10\_18（§\ref{sec:verify_dc_spectral}）では $\omega \in \{0.3, 0.5, 0.7\}$ すべてで
等密度（$\rho=1$）条件での収束を実験的に確認した（22--72 反復，tol=$10^{-8}$）．
$\omega = 0.83$ は Nyquist 減衰率 $0.992 \approx 1$ となり停滞する．
なお，$k = 3$ 反復で $\Ord{h^6}$ 精度に到達する（§\ref{sec:dc_iteration_accuracy}）理由は，
直接 LU（\texttt{spsolve}）による Step~2 の厳密求解により1反復あたりの有効減衰が大きいためである．

\subsubsection*{緩和係数 $\omega$ の役割}

式~\eqref{eq:dc_three_step} の Step~3 における $\omega$ は\textbf{直接緩和係数}であり，
Step~2 で直接 LU により厳密に求解した補正量 $\delta p^{(k+1)} = L_L^{-1} d^{(k)}$ に
スカラー倍 $\omega$ を適用する（式~\eqref{eq:dc_step2_lu}）．
収束条件 $\rho(M) < 1$ は固有値比 $\alpha_j \in [1.0, 2.4]$ の上限から
$\omega < 2/2.4 = 0.833$ として直接導かれる（式~\eqref{eq:dc_omega_max}）．

\subsubsection*{変密度場への拡張}

密度比 $\rho_l/\rho_g \gg 1$ の気液二相流では，
$L_H$ と $L_L$ の固有値比 $\alpha_j = \lambda_j/\mu_j$ が空間的に変動する．
特に：
\begin{itemize}
  \item 低密度相（気体）：$1/\rho$ が大きく $L_H$ の有効固有値が増大 → $\alpha_j$ が拡大
  \item 高密度相（液体）：$1/\rho$ が小さく有効固有値が低下 → $\alpha_j$ が縮小
\end{itemize}
この空間変動は $[\alpha_{\min}, \alpha_{\max}]$ の幅を広げ，
式~\eqref{eq:dc_rho_opt} の $\rho(M^*)$ を増大させる（収束が遅くなる）．
さらに，界面近傍では CCD の $\nabla\rho$ 評価（$\Ord{h^6}$）と
FD 前処理の $\nabla\rho$ 評価（$\Ord{h^2}$）の差異が「残差床」を形成し，
高密度比では DC+LU の収束が妨げられる（実験的確認: §\ref{sec:verify_dc_spectral}）．

\begin{tcolorbox}[resultbox, title={欠陥補正法の収束特性まとめ（DC + 直接 LU，Step 2 = \texttt{spsolve}）}]
\begin{center}
\PaperTableSetup
\begin{tabular}{lcc}
\toprule
条件 & $\rho_{\mathrm{Nyq}}(M)$ & 収束判定（tol $= 10^{-8}$） \\
\midrule
等密度，$\omega = 0.50$ & $0.20$ & 38 反復（$N=32$）収束 \\
等密度，$\omega = 0.70$ & $0.28$ & 22 反復（$N=32$）収束 \\
等密度，$\omega = 0.83\,(\omega_{\max})$ & $0.99$ & 停滞（不可） \\
変密度（$\rho_l/\rho_g = 2$，$N=64$） & $\approx 0.20$--$0.28$ & 39--74 反復 収束 \\
変密度（$\rho_l/\rho_g \ge 5$，$N=64$） & — & 界面 CCD/FD 不整合で停滞 \\
\bottomrule
\end{tabular}
\end{center}
収束には $\omega \in (0,\, 2/r_{\max}) = (0,\, 0.833)$ が必要である（実験検証: §\ref{sec:verify_dc_spectral}）．
変密度・鋭い界面では CCD（$\Ord{h^6}$）と FD（$\Ord{h^2}$）の $\nabla\rho$ 離散化の差異が
残差床 $\propto (\rho_l/\rho_g) \cdot h^2$ を生じるため，DC+LU の適用範囲は低密度比に限られる．
高密度比では一括 smoothed-Heaviside PPE の DC+LU に依存せず，
§~\ref{sec:split_ppe_framework} の分相 PPE + HFE 経路，または安定な FVM direct base solver を用いる．
\end{tcolorbox}
